% Edited and compiled online at https://letx.app
\documentclass[11pt,letterpaper]{article}

% Preamble
\usepackage[utf8]{inputenc}
\usepackage[T1]{fontenc}

% Typography
\usepackage{tgpagella} % TeX Gyre Pagella for professional body text
\usepackage[semibold]{sourcesanspro} % Source Sans Pro for sans-serif elements
\usepackage{microtype} % Improve justification and spacing

% Layout & Spacing
\usepackage{geometry}
\geometry{
    letterpaper,
    top=1.0in,
    bottom=1.0in,
    left=1.0in,
    right=1.0in
}
\usepackage{setspace}
\onehalfspacing % Clean, readable line spacing for academic statements

\usepackage{parskip}
\setlength{\parindent}{0pt}
\setlength{\parskip}{8pt plus 1pt minus 1pt}

% Colors
\usepackage{xcolor}
\definecolor{accent}{HTML}{1A365D}      % Deep navy for headings and highlights
\definecolor{textdark}{HTML}{2D3748}    % Charcoal for body text to reduce harshness
\definecolor{divider}{HTML}{CBD5E0}     % Subtle gray for lines

% Set default text color
\makeatletter
\newcommand{\globalcolor}[1]{%
  \color{#1}\global\let\default@color\current@color
}
\makeatother
\AtBeginDocument{\globalcolor{textdark}}

% Heading styling using titlesec
\usepackage{titlesec}
\titleformat{\section}
  {\sffamily\large\bfseries\color{accent}}
  {}{0pt}{}
  [\vspace{3pt}{\color{accent!30}\hrule height 1pt}\vspace{6pt}]
\titlespacing*{\section}{0pt}{2.5ex plus 0.5ex minus 0.2ex}{1.2ex plus 0.2ex}

% List styling using enumitem
\usepackage{enumitem}
\setlist[itemize]{
    leftmargin=1.5em,
    labelsep=0.5em,
    itemsep=5pt,
    topsep=3pt,
    parsep=0pt,
    partopsep=0pt
}

% Links and PDF metadata
\usepackage{hyperref}
\hypersetup{
    colorlinks=true,
    linkcolor=accent,
    urlcolor=accent,
    citecolor=accent,
    pdftitle={Diversity Statement - Dr. Marcus J. Carter},
    pdfauthor={Dr. Marcus J. Carter}
}

% Header and Footer
\usepackage{fancyhdr}
\pagestyle{fancy}
\fancyhf{}
\renewcommand{\headrulewidth}{0pt}
\renewcommand{\footrulewidth}{0pt}
\fancyfoot[L]{\sffamily\small\color{gray!80} Dr. Marcus J. Carter \quad$\cdot$\quad Diversity Statement}
\fancyfoot[R]{\sffamily\small\color{gray!80} Page \thepage}

\begin{document}

\begin{center}
    {\sffamily\Huge\bfseries\color{accent} Dr. Marcus J. Carter} \\
    \vspace{0.4em}
    {\sffamily\Large\color{textdark!80} Statement of Diversity, Equity, and Inclusion} \\
    \vspace{0.6em}
    {\color{divider}\rule{0.3\linewidth}{1.2pt}}
\end{center}
\vspace{1.5em}

\section{My Commitment to Diversity}

As a researcher and educator in Human-Computer Interaction (HCI), my career is centered on understanding how computing technology interacts with society. A core tenet of my professional philosophy is that academic excellence is inseparable from diversity, equity, and inclusion (DEI). Computing systems shape the modern world, yet the design of these systems and the classrooms in which they are taught have historically excluded marginalized voices. My commitment is to foster environments where students from all backgrounds---particularly women, underrepresented minorities, students with disabilities, and first-generation college students---can thrive. I view DEI not as an administrative mandate, but as a pedagogical and research imperative: diverse teams build more robust, creative, and ethical technologies, and equity in education is a fundamental right.

\section{Past Contributions \& Experience}

Throughout my academic career, I have actively worked to translate my commitment to diversity into concrete actions:

\begin{itemize}
    \item \textbf{Mentorship of Underrepresented Undergraduates:} Over the past four years, I have mentored eight undergraduate researchers, five of whom identify as women or minoritized individuals in engineering. Two of these students co-authored peer-reviewed publications and have since pursued graduate studies at top-tier institutions. I structured their research experiences with clear milestones and direct access to graduate mentors to demystify the research pathway.
    \item \textbf{K-12 Outreach in Underserved Communities:} I designed and co-led a bilingual summer computer science workshop, ``Code and Community,'' for middle school students in underfunded districts. By partnering with local community centers, we provided hands-on computing experience to 45 students, introducing them to computing concepts through the lens of local social issues.
    \item \textbf{Inclusive Curriculum Development:} As a graduate teaching assistant, I noticed that students without prior programming experience, often from low-resource high schools, faced higher drop rates in introductory courses. I developed a set of supplemental, scaffolded learning modules and organized weekly peer-led study groups. This initiative helped reduce the drop-fail-withdraw (DFW) rate in my sections by 18\% over three semesters.
\end{itemize}

\section{Inclusive Teaching \& Mentoring}

My teaching philosophy centers on culturally responsive pedagogy and Universal Design for Learning (UDL):

\begin{itemize}
    \item \textbf{Culturally Responsive Classrooms:} In my courses, I incorporate case studies and assignments that address systemic biases in technology, such as algorithmic fairness, racial biases in facial recognition, and the digital divide. By grounding technical concepts in real-world societal impact, I aim to create a learning environment where diverse students see their lives and concerns reflected in the syllabus.
    \item \textbf{Accessible Pedagogy:} Grounded in UDL principles, I structure my course materials to be accessible to all learning styles and abilities. I provide lecture transcripts, captioned videos, and alternative formats for project submissions (e.g., design portfolios, oral presentations, or traditional code repositories). I also implement a flexible grading framework that rewards growth and self-reflection.
    \item \textbf{Mentoring in the Lab:} I believe in proactive, structured mentorship. In my future research lab, I will implement Individual Development Plans (IDPs) for every student to align their personal career goals with their research tasks. I will establish clear communication guidelines to foster a culture of mutual respect and psychological safety, ensuring that all lab members feel valued and heard.
\end{itemize}

\section{Future Plans}

As a faculty member, I plan to build on these experiences to establish a sustained, positive impact:

\begin{itemize}
    \item \textbf{Sustaining Institutional DEI Goals:} I look forward to active collaboration with your institution's existing initiatives, such as the Women in Computer Science group and the campus-wide Office of Diversity and Inclusion. I plan to serve as a faculty advisor for student-led diversity organizations and participate in targeted recruiting events at regional conferences.
    \item \textbf{Community-Engaged Research:} My research lab will partner with local advocacy and accessibility groups to ensure that our computing research is informed by and beneficial to community needs. This participatory design approach will expose my students to the ethical and human dimensions of engineering.
    \item \textbf{Transfer Student Bridge Program:} Recognizing the systemic barriers faced by transfer students from community colleges, I aim to establish a winter bridge program that provides targeted academic advising, peer mentoring, and introductory research workshops to ease their transition into the computer science major.
\end{itemize}

\end{document}
